\subsection{Correctitud}
\label{sec:correctitud}

\begin{definicion}{Correctitud}{correctitud}
  Un algoritmo es \concept{correcto} si para toda instancia de entrada arroja la salida esperada, es decir, la salida que satisface la relación con la entrada definida por el problema.
\end{definicion}

Examinar la correctitud de un algoritmo es la primera etapa de su análisis y se le conoce como la \concept{validación} del algoritmo. Como la correctitud se debe verificar para todas las posibles instancias, usualmente exige una demostración matemática. Pero frecuentemente en la práctica basta con un argumento intuitivo o menos formal para convencerse de la validez de un algoritmo.

Sorprendentemente, incluso algunos algoritmos incorrectos pueden resultar útiles si producen una respuesta cercana a la esperada utilizando pocos recursos y no se conoce un algoritmo correcto o los algoritmos correctos conocidos son muy costosos. A estos algoritmos se les denomina \concept{algoritmos aproximados} y aunque son incorrectos se miden de acuerdo a qué tanto se acerca su solución a la esperada.

Un algoritmo puede ser \concept{determinístico}, si la respuesta que produce se puede determinar a partir de la entrada que recibe, o \concept{no determinístico}. En cualquiera de los dos casos, puede ser correcto o incorrecto, indistintamente.
